\scanpage{113}
\mspage{56}

Le lien avec la \deleted{définition}\inserted{description} (11) est évidemment que, si $\unsure{\xi|\zeta}_X$ provient d'un $p^t$ géom $\unsure{\xi|\zeta}_T$ de $T$, alors on aura des iso can

(20)\quad
\diagram{\begin{tikzcd}[column sep=large,row sep=large]
\mathcal E_{X/Y}^{\unsure{\xi|\zeta}_T}\arrow[r,phantom,"\simeq"]\arrow[d,phantom,"\rotatebox{90}{$\simeq$}" description] & \mathcal E_{X/Y}^T\\
\pi_1(X_{\unsure{\xi|\zeta}_Y},\unsure{\xi|\zeta}_T) &
\end{tikzcd}}

\deleted{\unsure{donc}} --- les $\pi_1(X_{\unsure{\xi|\zeta}_Y},\unsure{\xi|\zeta}_T)$ des fibres géom, relatifs aux divers pts géom de $\mathcal E_{X/Y}^T$, sont canoniquement iso entre eux.

( Il faut bien faire attention que $\mathcal E_{X/Y}$ \inserted{($\unsure{Z|T}$ non précisé)} est une classe d'iso. de groupes profinis sans plus, \uline{pas} un groupe extérieur. L'arbitraire dans le choix d'un iso entre $\mathcal E_{X/Y}^T$ et $\mathcal E_{X/Y}^{T'}$ provient du choix

\margin{112}
